% Topic 1.2.13 — Goodness-of-Fit: Kolmogorov and Chi-Squared Tests.

\section{Критерии согласия Колмогорова и Пирсона}
\label{sec:1.2.13-goodness-of-fit}

\begin{ticket}
Эмпирическая функция распределения. Статистическая гипотеза, выборка, критическая область гипотезы, уровень значимости. Теоремы о свойствах критериев согласия Колмогорова и Пирсона (хи-квадрат). Теорема Крамера (параметрический хи-квадрат). Применение критериев для проверки согласия результатов опыта с теоретическими гипотезами. Использование статистических таблиц.
\end{ticket}

\subsection*{Теория}

\emph{Проверка согласия} (\emph{goodness of fit}) ставит вопрос: правдоподобно ли, что наблюдаемые данные получены из некоторого заданного распределения. На практике основная масса учебных задач закрывается двумя классическими непараметрическими критериями: критерием Колмогорова--Смирнова (для непрерывной~$F$) и критерием Пирсона хи-квадрат (для дискретных или сгруппированных данных). Оба опираются на асимптотическое распределение явной статистики, квантили которого затабулированы.

\medskip
\noindent\textbf{Эмпирическая функция распределения.}

\begin{definition}[Эмпирическая функция распределения]
\label{def:empirical-cdf}
Пусть $X_1, \dots, X_n$~--- выборка (независимые одинаково распределённые наблюдения) с общей функцией распределения~$F$. \emph{Эмпирической функцией распределения} называется
\[
  F_n(x) \;=\; \frac{1}{n} \sum_{i=1}^{n} \mathbf{1}_{\{X_i \leq x\}}, \qquad x \in \R.
\]
При фиксированном~$x$ имеем $n F_n(x) \sim \mathrm{Bin}(n, F(x))$, поэтому $\E F_n(x) = F(x)$ и $\Var F_n(x) = F(x)(1 - F(x))/n$.
\end{definition}

\begin{theorem}[Гливенко--Кантелли]
\label{thm:glivenko-cantelli}
Для выборки независимых одинаково распределённых величин с непрерывной функцией распределения~$F$
\[
  \sup_{x \in \R}\, |F_n(x) - F(x)| \;\xrightarrow{\text{п.н.}}\; 0 \qquad (n \to \infty).
\]
\end{theorem}
\proofof{glivenko-cantelli}

Теорема Гливенко--Кантелли утверждает, что эмпирическая функция распределения~--- равномерно состоятельная оценка истинной. Статистики критериев согласия измеряют скорость этой сходимости.

\medskip
\noindent\textbf{Критерий Колмогорова--Смирнова.}

\begin{theorem}[Колмогоров]
\label{thm:kolmogorov-test}
При простой нулевой гипотезе $H_0: X_i \sim F$ с непрерывной~$F$ статистика
\[
  D_n \;=\; \sqrt n\, \sup_{x \in \R} |F_n(x) - F(x)|
\]
сходится по распределению к \emph{распределению Колмогорова}~$K$ с функцией распределения
\[
  Q(t) = 1 - 2 \sum_{k=1}^{\infty} (-1)^{k-1} e^{-2 k^2 t^2}, \qquad t \geq 0.
\]
Это распределение \emph{не зависит от~$F$} (свободно от распределения).
\end{theorem}
\proofof{kolmogorov-test}

\emph{Критерий согласия Колмогорова} размера~$\alpha$ отвергает~$H_0$, когда $D_n$ превышает верхнюю $\alpha$-квантиль~$K_\alpha$ распределения~$K$ (берётся из таблиц; например, $K_{0{,}05} \approx 1{,}358$). Критерий асимптотически свободен от распределения и состоятелен против любой фиксированной непрерывной альтернативы $G \neq F$ (поскольку $\sup |F_n - F|$ при альтернативе сходится к $\sup |G - F| > 0$, тогда как при нулевой гипотезе обращается в нуль).

\medskip
\noindent\textbf{Критерий Пирсона хи-квадрат.}

\begin{definition}[Мультиномиальная схема]
\label{def:multinomial}
Разобьём носитель~$F$ на~$r$ непересекающихся ячеек $A_1, \dots, A_r$. Положим $p_k = \Prob_F(X \in A_k)$ (вероятности ячеек при~$F$) и $n_k = \#\{i : X_i \in A_k\}$ (наблюдаемые частоты). Тогда $(n_1, \dots, n_r) \sim \mathrm{Multinomial}(n; p_1, \dots, p_r)$, причём $\E n_k = n p_k$ и $\Var n_k = n p_k (1 - p_k)$.
\end{definition}

\begin{theorem}[Пирсон, хи-квадрат]
\label{thm:pearson-chi2}
При нулевой гипотезе $H_0 : (n_1, \dots, n_r) \sim \mathrm{Multinomial}(n; p_1, \dots, p_r)$ со всеми $p_k > 0$ статистика
\[
  \chi^2 \;=\; \sum_{k=1}^{r} \frac{(n_k - n p_k)^2}{n p_k}
\]
сходится по распределению к $\chi^2_{r - 1}$ при $n \to \infty$.
\end{theorem}
\proofof{pearson-chi2}

\emph{Критерий Пирсона хи-квадрат} размера~$\alpha$ отвергает~$H_0$ при $\chi^2 > \chi^2_{\alpha, r - 1}$ (верхняя $\alpha$-квантиль распределения $\chi^2_{r-1}$, по таблицам). Стандартное практическое требование~--- $n p_k \geq 5$ в каждой ячейке, чтобы асимптотическая аппроксимация была надёжной.

\begin{theorem}[Крамер: параметрический хи-квадрат]
\label{thm:cramer-parametric}
Пусть~$H_0$ сложна и имеет вид ``$F \in \{F_\theta : \theta \in \Theta\}$'' при $\Theta \subseteq \R^s$, а $\hat \theta$~--- эффективная (например, оценка максимального правдоподобия) оценка параметра~$\theta$. Зададим оценочные вероятности ячеек $\hat p_k = \Prob_{\hat \theta}(X \in A_k)$. Тогда при $H_0$
\[
  \chi^2 \;=\; \sum_{k=1}^{r} \frac{(n_k - n \hat p_k)^2}{n \hat p_k} \;\xrightarrow{d}\; \chi^2_{r - 1 - s}.
\]
\end{theorem}
\proofof{cramer-parametric}

\begin{remark}
Каждый параметр, \emph{оценённый по той же выборке}, отнимает одну степень свободы у предельного $\chi^2$. Тот же приём лежит в основе критерия независимости в таблицах сопряжённости ($r$ строк, $c$ столбцов; число степеней свободы $(r - 1)(c - 1)$). Неасимптотический вариант критерия~$\chi^2$ для конечных выборок изложен в~\cite[гл.~VIII]{Sevastyanov1982}; статистические таблицы для распределений $K$, $\chi^2$, $t$ и $F$ традиционно индексируются по верхнехвостовой вероятности.
\end{remark}

\subsection*{Доказательства}

\begin{proof}[\Cref{thm:glivenko-cantelli}, набросок]
\proofanchor{glivenko-cantelli}
Применим усиленный закон больших чисел (\cref{thm:slln}) при каждом фиксированном~$x$ к индикаторам $\mathbf{1}_{\{X_i \leq x\}}$, независимым и одинаково распределённым, с математическим ожиданием~$F(x)$: получим $F_n(x) \to F(x)$ почти наверное. Усиление до равномерной сходимости использует (i) монотонность функций~$F_n$ и~$F$ и (ii) непрерывность справа и существование пределов слева у~$F$. Возьмём конечную сетку $-\infty = x_0 < x_1 < \cdots < x_m = +\infty$ так, чтобы $F(x_j) - F(x_{j-1}^{-}) < \varepsilon$ при всех~$j$ (это возможно, выбрав~$x_j$ как соответствующие квантили~$F$). На событии полной меры, где $F_n(x_j) \to F(x_j)$ и $F_n(x_j^-) \to F(x_j^-)$ при всех~$j$ (пересечение счётного числа событий полной вероятности также имеет полную вероятность), монотонность зажимает $F_n(x)$ между $F_n(x_{j-1})$ и $F_n(x_j)$ при $x \in [x_{j-1}, x_j)$, так что $\sup |F_n - F|$ на всей прямой не превосходит~$\varepsilon$ плюс стремящаяся к нулю поточечная погрешность. Устремляя $\varepsilon \to 0$, получаем требуемое. Полное доказательство~--- в~\cite[гл.~II, \S~10]{Shiryaev1989}.
\end{proof}

\begin{proof}[\Cref{thm:kolmogorov-test}, идея]
\proofanchor{kolmogorov-test}
Преобразование вероятностной шкалы $U_i = F(X_i) \sim \mathrm{Uniform}(0, 1)$ освобождает статистику от зависимости от~$F$: $D_n$ по выборке~$(X_i)$ при законе~$F$ совпадает с $D_n$ по выборке~$(U_i)$ при равномерном законе. Эмпирический процесс $\sqrt n (F_n - F)$ по теореме Донскера (функциональная ЦПТ) сходится к \emph{броуновскому мосту}~$B^0_t$ на отрезке $[0, 1]$, супремум модуля которого имеет распределение $\sup_t |B^0_t| \sim K$. Вероятность $\Prob(\sup_t |B^0_t| \leq t)$ вычисляется по принципу отражения для броуновского движения, что и приводит к знакочередующейся сумме~$Q(t)$. См.~\cite[гл.~XII]{Shiryaev1989} для подхода через функциональную ЦПТ или элементарное комбинаторное доказательство в~\cite[гл.~IX]{Sevastyanov1982}.
\end{proof}

\begin{proof}[\Cref{thm:pearson-chi2}, набросок]
\proofanchor{pearson-chi2}
Введём нормированные частоты $Y_k = (n_k - n p_k)/\sqrt{n p_k}$. По многомерной ЦПТ, применённой к мультиномиальному вектору, случайный вектор $(Y_1, \dots, Y_r)$ сходится по распределению к центрированному гауссовскому $\vect{Y} \sim \mathcal{N}(\vect{0}, \mat{I}_r - \sqrt{\vect{p}}\,\sqrt{\vect{p}}\T)$, где $\sqrt{\vect{p}} = (\sqrt{p_1}, \dots, \sqrt{p_r})\T$~--- единичный вектор ($\|\sqrt{\vect{p}}\|^2 = \sum p_k = 1$). Ковариационная матрица $\mat{I} - \sqrt{\vect{p}}\,\sqrt{\vect{p}}\T$ есть проектор на $(r - 1)$-мерное гиперподпространство $\sqrt{\vect{p}}^{\perp}$~--- симметричная идемпотентная матрица ранга $r - 1$.

Статистика Пирсона $\chi^2 = \sum Y_k^2 = \|\vect{Y}\|^2$ поэтому сходится по распределению к квадрату нормы стандартного гауссовского вектора в подпространстве $\sqrt{\vect{p}}^{\perp}$, то есть к $\chi^2_{r - 1}$ согласно~\cref{def:chi-squared} (любой ортонормированный базис $(r-1)$-мерного подпространства даёт $r - 1$ независимых стандартных нормальных).
\end{proof}

\begin{proof}[\Cref{thm:cramer-parametric}, идея]
\proofanchor{cramer-parametric}
В статистике Пирсона заменим истинные вероятности ячеек $p_k(\theta_0)$ на их оценки $p_k(\hat \theta)$. Разлагая в ряд Тейлора вокруг истинного значения~$\theta_0$, получаем, что $\sqrt n (\hat \theta - \theta_0)$ ``съедает''~$s$ степеней свободы у предельного гауссовского вектора: предел живёт в ортогональном дополнении одновременно и к $\sqrt{\vect{p}(\theta_0)}$, и к $s$-мерному подпространству, натянутому на нормированные градиенты $\sqrt{n}\, \partial p_k/\partial \theta_j$ в точке~$\theta_0$. Получившееся предельное распределение~--- $\chi^2_{r - 1 - s}$. Условие эффективности~$\hat \theta$ (асимптотическая дисперсия совпадает с границей Крамера--Рао) и обеспечивает работоспособность проективно-геометрического аргумента; иначе предел оказался бы нецентральным или взвешенным $\chi^2$, а не чистым~$\chi^2$. Подробный вывод~--- в~\cite[гл.~IX, \S~3]{Sevastyanov1982}.
\end{proof}

\subsection*{Примеры}

\begin{example}[Правильная ли кость, критерий Пирсона]
\label{ex:fair-die-test}
Кость бросают $n = 600$ раз, наблюдаемые частоты $(n_1, \dots, n_6) = (95, 105, 100, 110, 90, 100)$. Проверим $H_0$~--- кость правильна, $p_k = 1/6$ при всех~$k$. Ожидаемые частоты $n p_k = 100$. Статистика Пирсона:
\[
  \chi^2 = \frac{(-5)^2 + 5^2 + 0 + 10^2 + (-10)^2 + 0}{100} = \frac{25 + 25 + 100 + 100}{100} = 2{,}5.
\]
Сравним с верхней $0{,}05$-квантилью $\chi^2_5$, равной $\chi^2_{0{,}05, 5} \approx 11{,}07$. Поскольку $2{,}5 \ll 11{,}07$, нулевую гипотезу на уровне~$0{,}05$ не отвергаем~--- данные согласуются с правильной костью.
\end{example}

\begin{example}[Колмогоров--Смирнов против равномерного]
\label{ex:ks-uniform}
Имеется выборка размера $n = 100$, проверяемая на согласие с распределением $\mathrm{Uniform}(0, 1)$; наибольшее отклонение $\sup_x |F_n(x) - x|$ оказалось равным~$0{,}10$. Тогда $D_n = \sqrt{100} \cdot 0{,}10 = 1{,}00$. Верхняя $0{,}05$-квантиль распределения Колмогорова~--- $K_{0{,}05} \approx 1{,}358$. Поскольку $1{,}00 < 1{,}358$, нулевую гипотезу на уровне~$0{,}05$ не отвергаем. Большее отклонение в~$0{,}15$ дало бы $D_n = 1{,}5 > 1{,}358$ и привело бы к отказу от нулевой гипотезы.
\end{example}

\begin{problem}
\label{prob:cramer-parametric-poisson}
Завод производит лампочки, время жизни которых тестируется партиями по~$n = 200$ штук. Предполагается, что время жизни распределено по закону $\mathrm{Exp}(\lambda)$ при неизвестном~$\lambda$. Времена жизни сгруппированы в~$r = 5$ ячеек; наблюдаемые и (оцененные) ожидаемые частоты:
$\vect{n} = (52, 38, 39, 35, 36)$ и $n \hat{\vect{p}} = (50, 42, 35, 31, 42)$, где $\hat \lambda$~--- оценка максимального правдоподобия. Найдите статистику хи-квадрат и сделайте вывод на уровне $\alpha = 0{,}05$.
\end{problem}

\begin{proof}[Решение]
Подставляем в статистику Пирсона:
\[
  \chi^2 = \frac{4}{50} + \frac{16}{42} + \frac{16}{35} + \frac{16}{31} + \frac{36}{42} \approx 0{,}080 + 0{,}381 + 0{,}457 + 0{,}516 + 0{,}857 \approx 2{,}29.
\]
По теореме Крамера (\cref{thm:cramer-parametric}), поскольку из выборки оценивался $s = 1$ параметр (a именно~$\lambda$), предельное распределение~--- $\chi^2_{r - 1 - s} = \chi^2_3$. Верхняя $0{,}05$-квантиль $\chi^2_3$ равна~$7{,}815$. Поскольку $2{,}29 \ll 7{,}815$, экспоненциальную нулевую гипотезу на уровне~$0{,}05$ не отвергаем~--- данные согласуются с моделью экспоненциального времени жизни.
\end{proof}
